\documentclass[a4paper, 14pt]{extarticle}

% ============ Пакеты ============
\usepackage[utf8]{inputenc}
\usepackage[T2A]{fontenc}
\usepackage[russian]{babel}
\usepackage{geometry}
\geometry{left=2.5cm, right=1.5cm, top=2cm, bottom=2cm}

\usepackage{amsmath, amssymb, amsthm}
\usepackage{graphicx}
\usepackage{float}
\usepackage{caption}
\usepackage{subcaption}
\usepackage{booktabs}
\usepackage{listings}
\usepackage{xcolor}
\usepackage{hyperref}
\usepackage{indentfirst}
\usepackage{setspace}
\usepackage{enumitem}
\usepackage{titlesec}

\onehalfspacing

\hypersetup{
    colorlinks=true,
    linkcolor=blue!60!black,
    citecolor=blue!60!black,
    urlcolor=blue!60!black
}

% ============ Настройки листингов ============
\definecolor{codegreen}{rgb}{0,0.5,0}
\definecolor{codegray}{rgb}{0.5,0.5,0.5}
\definecolor{codepurple}{rgb}{0.58,0,0.82}
\definecolor{backcolour}{rgb}{0.97,0.97,0.97}

\lstdefinestyle{python}{
    backgroundcolor=\color{backcolour},
    commentstyle=\color{codegreen},
    keywordstyle=\color{blue},
    numberstyle=\tiny\color{codegray},
    stringstyle=\color{codepurple},
    basicstyle=\ttfamily\small,
    breaklines=true,
    numbers=left,
    numbersep=5pt,
    showstringspaces=false,
    tabsize=4,
    language=Python,
    frame=single,
    rulecolor=\color{codegray}
}

% ============ Теоремы ============
\newtheorem{definition}{Определение}
\newtheorem{theorem}{Теорема}
\newtheorem{remark}{Замечание}

% ============ Начало документа ============
\begin{document}

% ============ Титульная страница ============
\begin{titlepage}
    \centering
    \vspace*{1cm}

    {\large Университет}\\[0.3cm]
    {\large Факультет}\\[0.3cm]
    {\large Кафедра}\\[2cm]

    {\LARGE\bfseries Лабораторная работа №1}\\[0.5cm]
    {\Large Методы оптимизации}\\[0.3cm]
    {\large Курс <<Продвинутые методы оптимизации>>}\\[3cm]

    \begin{flushright}
        \large
        Выполнил: студент группы \\
        ФИО \\[0.5cm]
        Преподаватель: \\
        ФИО
    \end{flushright}

    \vfill
    {\large 2026}
\end{titlepage}

% ============ Оглавление ============
\tableofcontents
\newpage

% ============ Введение ============
\section{Введение}

Целью данной лабораторной работы является реализация типа данных <<конструктивное число>>, набора целевых функций (<<чёрных ящиков>>) и методов оптимизации различных порядков, а также проведение исследования сходимости методов и поведения погрешности конструктивных чисел в процессе оптимизации.

В рамках работы реализованы:
\begin{itemize}
    \item \textbf{Конструктивное число} --- тип данных на основе интервальной арифметики;
    \item \textbf{Чёрные ящики} --- три целевые функции с унифицированным интерфейсом;
    \item \textbf{Метод 0-го порядка} --- метод деформируемого многогранника (Нелдера--Мида);
    \item \textbf{Метод 1-го порядка} --- градиентный спуск;
    \item \textbf{Метод 2-го порядка} --- метод Ньютона.
\end{itemize}

% ================================================================
%                       ТЕОРЕТИЧЕСКАЯ ЧАСТЬ
% ================================================================
\section{Теоретическая часть}

% -------- 2.1 Конструктивное число --------
\subsection{Конструктивное число}

\begin{definition}
\textbf{Конструктивное число} --- пара $(a, b) \in \mathbb{Q}^2$, $a \le b$, задающая интервал $[a, b]$, в котором находится представляемое действительное число. Величина $\varepsilon = b - a$ характеризует точность представления.
\end{definition}

Конструктивное число может быть создано двумя способами:
\begin{enumerate}
    \item Из числа $x \in \mathbb{R}$ и точности $\varepsilon > 0$:
    \[
        a = x - \frac{\varepsilon}{2}, \quad b = x + \frac{\varepsilon}{2}.
    \]
    \item Из пары рациональных чисел $a, b \in \mathbb{Q}$ напрямую.
\end{enumerate}

Получение действительного числа из конструктивного производится через параметр $\alpha \in [0, 1]$:
\[
    x(\alpha) = a \cdot (1 - \alpha) + b \cdot \alpha.
\]
При $\alpha = 0$ получаем $a$, при $\alpha = 1$ получаем $b$, при $\alpha = 0.5$ --- середину интервала.

\subsubsection{Интервальная арифметика}

Арифметические операции над конструктивными числами определяются правилами интервальной арифметики. Пусть $\mathbf{x} = [a_1, b_1]$ и $\mathbf{y} = [a_2, b_2]$:

\begin{align}
    \mathbf{x} + \mathbf{y} &= [a_1 + a_2,\; b_1 + b_2], \label{eq:add}\\
    \mathbf{x} - \mathbf{y} &= [a_1 - b_2,\; b_1 - a_2], \label{eq:sub}\\
    \mathbf{x} \cdot \mathbf{y} &= \bigl[\min\{a_1 a_2, a_1 b_2, b_1 a_2, b_1 b_2\},\; \max\{a_1 a_2, a_1 b_2, b_1 a_2, b_1 b_2\}\bigr], \label{eq:mul}\\
    \mathbf{x} / \mathbf{y} &= \mathbf{x} \cdot [1/b_2,\; 1/a_2], \quad \text{при } 0 \notin [a_2, b_2]. \label{eq:div}
\end{align}

\begin{remark}
Из формул~\eqref{eq:add}--\eqref{eq:div} видно, что ширина интервала (погрешность $\varepsilon$) монотонно возрастает при каждой операции. Это является фундаментальным свойством интервальной арифметики и ключевым предметом исследования данной работы.
\end{remark}

Сравнение интервалов:
\[
    \mathbf{x} < \mathbf{y} \iff b_1 < a_2, \qquad
    \mathbf{x} > \mathbf{y} \iff a_1 > b_2.
\]

% -------- 2.2 Целевые функции --------
\subsection{Целевые функции}

\subsubsection{Квадратичная функция}

Общий вид квадратичной функции $n$ переменных:
\begin{equation}\label{eq:quadratic}
    f(\mathbf{x}) = \frac{1}{2}\mathbf{x}^T A \mathbf{x} + \mathbf{b}^T \mathbf{x} + c,
\end{equation}
где $A \in \mathbb{R}^{n \times n}$ --- симметричная положительно определённая матрица, $\mathbf{b} \in \mathbb{R}^n$, $c \in \mathbb{R}$.

Градиент и гессиан:
\[
    \nabla f(\mathbf{x}) = A\mathbf{x} + \mathbf{b}, \qquad H_f = A.
\]

Минимум достигается в точке $\mathbf{x}^* = -A^{-1}\mathbf{b}$.

\begin{definition}
\textbf{Число обусловленности} матрицы $A$:
\[
    \kappa(A) = \frac{\lambda_{\max}(A)}{\lambda_{\min}(A)},
\]
где $\lambda_{\max}$, $\lambda_{\min}$ --- наибольшее и наименьшее собственные значения $A$.
\end{definition}

Число обусловленности определяет скорость сходимости градиентных методов. При $\kappa \approx 1$ линии уровня близки к сферам, и методы сходятся быстро. При большом $\kappa$ линии уровня вытянуты, и сходимость замедляется.

В работе используются:
\begin{itemize}
    \item \textbf{Функция 1}: $n=6$, собственные значения $\{0.95, 0.98, 1.0, 1.02, 1.05, 1.1\}$, $\kappa \approx 1.16$.
    \item \textbf{Функция 2}: $n=4$, собственные значения $\{1, 10, 50, 100\}$, $\kappa = 100$.
\end{itemize}

Матрицы $A$ конструируются с заданным спектром через случайное ортогональное преобразование: $A = Q \Lambda Q^T$, где $Q$ --- ортогональная матрица, $\Lambda = \mathrm{diag}(\lambda_1, \ldots, \lambda_n)$.

\subsubsection{Функция Розенброка}

Обобщённая функция Розенброка для $n$ переменных:
\begin{equation}\label{eq:rosenbrock}
    f(\mathbf{x}) = \sum_{i=1}^{n-1} \left[ 100\,(x_{i+1} - x_i^2)^2 + (1 - x_i)^2 \right].
\end{equation}

В работе используется при $n = 3$. Глобальный минимум $f(\mathbf{x}^*) = 0$ достигается в $\mathbf{x}^* = (1, 1, 1)^T$.

Градиент:
\begin{align}
    \frac{\partial f}{\partial x_i} &= -400\,x_i\,(x_{i+1} - x_i^2) + 2\,(x_i - 1), \quad i < n, \\
    \frac{\partial f}{\partial x_n} &= 200\,(x_n - x_{n-1}^2).
\end{align}

Гессиан является трёхдиагональной матрицей:
\[
    H_{ii} = -400(x_{i+1} - 3x_i^2) + 2, \quad
    H_{i,i+1} = H_{i+1,i} = -400\,x_i, \quad
    H_{nn} = 200.
\]

Функция Розенброка является классическим тестом для методов оптимизации: она имеет узкую изогнутую <<долину>>, по которой методам сложно двигаться.

% -------- 2.3 Методы оптимизации --------
\subsection{Методы оптимизации}

\subsubsection{Метод деформируемого многогранника (Нелдера--Мида)}

Метод Нелдера--Мида --- метод \textbf{0-го порядка} (не использует производные). Работает с симплексом из $n+1$ вершин в $\mathbb{R}^n$.

\begin{enumerate}
    \item \textbf{Инициализация}: строится начальный симплекс из $n+1$ точек.
    \item \textbf{Сортировка}: вершины упорядочиваются по значению функции:
    \[
        f(\mathbf{x}_1) \le f(\mathbf{x}_2) \le \ldots \le f(\mathbf{x}_{n+1}).
    \]
    \item \textbf{Центроид}: $\mathbf{x}_c = \frac{1}{n}\sum_{i=1}^{n} \mathbf{x}_i$ (без худшей точки).
    \item \textbf{Отражение}: $\mathbf{x}_r = \mathbf{x}_c + \alpha(\mathbf{x}_c - \mathbf{x}_{n+1})$, $\alpha = 1$.
    \item \textbf{Растяжение}: если $f(\mathbf{x}_r) < f(\mathbf{x}_1)$, то $\mathbf{x}_e = \mathbf{x}_c + \gamma(\mathbf{x}_r - \mathbf{x}_c)$, $\gamma = 2$.
    \item \textbf{Сжатие}: если $f(\mathbf{x}_r) \ge f(\mathbf{x}_n)$, то $\mathbf{x}_s = \mathbf{x}_c + \rho(\mathbf{x}_{n+1} - \mathbf{x}_c)$, $\rho = 0.5$.
    \item \textbf{Глобальное сжатие}: $\mathbf{x}_i \leftarrow \mathbf{x}_1 + \sigma(\mathbf{x}_i - \mathbf{x}_1)$, $\sigma = 0.5$.
\end{enumerate}

Критерий остановки --- стандартное отклонение значений $f$ в вершинах симплекса меньше заданного $\varepsilon_{\mathrm{tol}}$.

\subsubsection{Градиентный спуск}

Метод \textbf{1-го порядка}. Итеративная формула:
\begin{equation}\label{eq:gd}
    \mathbf{x}_{k+1} = \mathbf{x}_k - \eta \nabla f(\mathbf{x}_k),
\end{equation}
где $\eta > 0$ --- скорость обучения (learning rate).

\begin{theorem}[Скорость сходимости]
Для сильно выпуклой функции с числом обусловленности $\kappa$ при оптимальном выборе $\eta = \frac{2}{\lambda_{\max} + \lambda_{\min}}$:
\[
    \|\mathbf{x}_k - \mathbf{x}^*\| \le \left(\frac{\kappa - 1}{\kappa + 1}\right)^k \|\mathbf{x}_0 - \mathbf{x}^*\|.
\]
\end{theorem}

Таким образом, при $\kappa \approx 1$ сходимость быстрая (геометрическая прогрессия с малым знаменателем), а при $\kappa \gg 1$ --- медленная.

Выбор $\eta$ критически важен:
\begin{itemize}
    \item Слишком большой $\eta$ $\Rightarrow$ расходимость;
    \item Слишком малый $\eta$ $\Rightarrow$ медленная сходимость;
    \item Оптимальный $\eta \in (0, 2/\lambda_{\max})$.
\end{itemize}

\subsubsection{Метод Ньютона}

Метод \textbf{2-го порядка}. Итеративная формула:
\begin{equation}\label{eq:newton}
    \mathbf{x}_{k+1} = \mathbf{x}_k - [H_f(\mathbf{x}_k)]^{-1} \nabla f(\mathbf{x}_k),
\end{equation}
где $H_f(\mathbf{x}_k)$ --- гессиан функции в точке $\mathbf{x}_k$.

Для квадратичных функций метод Ньютона сходится за \textbf{одну итерацию} (точное решение СЛАУ).

Для общих функций вблизи минимума наблюдается \textbf{квадратичная сходимость}:
\[
    \|\mathbf{x}_{k+1} - \mathbf{x}^*\| \le C \|\mathbf{x}_k - \mathbf{x}^*\|^2.
\]

Достоинства:
\begin{itemize}
    \item Не зависит от числа обусловленности;
    \item Квадратичная скорость сходимости вблизи минимума.
\end{itemize}

Недостатки:
\begin{itemize}
    \item Требует вычисления и обращения гессиана --- $O(n^3)$ на итерацию;
    \item Может расходиться, если гессиан не положительно определён.
\end{itemize}


% ================================================================
%                      РЕАЛИЗАЦИЯ
% ================================================================
\section{Описание реализации}

\subsection{Структура проекта}

\begin{verbatim}
src/
    __init__.py
    ConstructiveNumber.py   -- конструктивное число
    BlackBox.py             -- чёрные ящики функций
    Optimizers.py           -- методы оптимизации
test.ipynb                  -- тесты
research.ipynb              -- исследования и графики
report.tex                  -- данный отчёт
\end{verbatim}

\subsection{Конструктивное число (\texttt{ConstructiveNumber})}

Класс хранит пару \texttt{(a, b)} и реализует:

\begin{lstlisting}[style=python, caption={Создание конструктивного числа}]
# Из пары (a, b)
c = ConstructiveNumber(1.0, 2.0)

# Из числа x и точности epsilon
c = ConstructiveNumber(5.0, epsilon=0.01)
\end{lstlisting}

Перегружены операторы Python (\texttt{\_\_add\_\_}, \texttt{\_\_sub\_\_}, \texttt{\_\_mul\_\_}, \texttt{\_\_truediv\_\_}, \texttt{\_\_pow\_\_}, \texttt{\_\_lt\_\_}, \texttt{\_\_gt\_\_} и др.), что обеспечивает нативный синтаксис: \texttt{a + b}, \texttt{a * b}, \texttt{a < b}.

Каждая операция поддерживает как пару \texttt{CN + CN}, так и \texttt{CN + float} (и \texttt{float + CN} через \texttt{\_\_radd\_\_} и аналоги).

Возведение в квадрат оптимизировано: для $[a, b]$ при $a \ge 0$ результат $[a^2, b^2]$, при $b \le 0$ --- $[b^2, a^2]$, иначе --- $[0, \max(a^2, b^2)]$.

\subsection{Чёрные ящики (\texttt{BlackBox})}

Базовый класс \texttt{BlackBox} определяет интерфейс:
\begin{itemize}
    \item \texttt{func(x)} --- значение функции;
    \item \texttt{grad(x)} --- градиент;
    \item \texttt{hess(x)} --- гессиан;
    \item \texttt{dim} --- размерность;
    \item \texttt{name} --- название;
    \item \texttt{minimum} --- известная точка минимума;
    \item Счётчики: \texttt{func\_calls}, \texttt{grad\_calls}, \texttt{hess\_calls}.
\end{itemize}

Все три функции (\texttt{QuadraticGoodCond}, \texttt{QuadraticBadCond}, \texttt{Rosenbrock}) при вызове проверяют тип входных данных: если \texttt{x[0]} является \texttt{ConstructiveNumber}, вычисления проводятся в интервальной арифметике.

\subsection{Методы оптимизации (\texttt{Optimizers})}

Каждый метод возвращает пару \texttt{(x\_result, history)}, где \texttt{history} содержит:
\begin{itemize}
    \item \texttt{x} --- траектория точек;
    \item \texttt{f} --- значения функции;
    \item \texttt{grad\_norm} --- нормы градиента (для GD и Newton);
    \item \texttt{epsilon} --- динамика $\varepsilon$ при использовании CN;
    \item \texttt{time} --- время каждой итерации.
\end{itemize}

При флаге \texttt{use\_cn=True} входной вектор преобразуется в массив \texttt{ConstructiveNumber} с заданной начальной точностью \texttt{cn\_epsilon}, и все вычисления идут в интервальной арифметике.


% ================================================================
%                      ИССЛЕДОВАНИЯ
% ================================================================
\section{Исследования}

\subsection{Подбор learning rate}

Для каждой функции был проведён перебор значений $\eta$ (рис.~\ref{fig:lr}).

Выбранные оптимальные значения:
\begin{itemize}
    \item Квадратичная 6D ($\kappa \approx 1.16$): $\eta = 1.5$;
    \item Квадратичная 4D ($\kappa = 100$): $\eta = 0.018$;
    \item Розенброк 3D: $\eta = 0.003$.
\end{itemize}

\begin{figure}[H]
    \centering
    \includegraphics[width=\textwidth]{img/lr_selection.png}
    \caption{Подбор learning rate для градиентного спуска}
    \label{fig:lr}
\end{figure}

Видно, что для хорошо обусловленной задачи допустим широкий диапазон $\eta$, тогда как для плохо обусловленной --- допустимый диапазон крайне узок ($\eta < 2/\lambda_{\max} = 0.02$).

\subsection{Сравнение методов оптимизации}

\subsubsection{Сводная таблица}

Результаты запуска всех методов на всех функциях из начальной точки $\mathbf{x}_0 = \mathbf{0}$ (для квадратичных) и $\mathbf{x}_0 = (-1, -1, -1)^T$ (для Розенброка):

\begin{table}[H]
\centering
\caption{Сравнение методов оптимизации}
\label{tab:comparison}
\begin{tabular}{llrrrr}
\toprule
Функция & Метод & Итерации & Выз.~$f$ & Выз.~$\nabla f$ & $\|\mathbf{x} - \mathbf{x}^*\|$ \\
\midrule
Квадр.~6D ($\kappa\approx1$) & Нелдер--Мид & $\sim 200$ & $\sim 400$ & 0 & $\sim 10^{-5}$ \\
                               & Град. спуск & $\sim 30$ & $\sim 30$ & $\sim 30$ & $\sim 10^{-9}$ \\
                               & Ньютон & 2 & 2 & 2 & $< 10^{-15}$ \\
\midrule
Квадр.~4D ($\kappa=100$) & Нелдер--Мид & $\sim 300$ & $\sim 600$ & 0 & $\sim 10^{-5}$ \\
                           & Град. спуск & $\sim 5000$ & $\sim 5000$ & $\sim 5000$ & $\sim 10^{-3}$ \\
                           & Ньютон & 2 & 2 & 2 & $< 10^{-15}$ \\
\midrule
Розенброк 3D & Нелдер--Мид & $\sim 300$ & $\sim 600$ & 0 & $\sim 10^{-4}$ \\
              & Град. спуск & $\sim 5000$ & $\sim 5000$ & $\sim 5000$ & $\sim 10^{-1}$ \\
              & Ньютон & $\sim 10$ & $\sim 10$ & $\sim 10$ & $< 10^{-10}$ \\
\bottomrule
\end{tabular}
\end{table}

\begin{remark}
Конкретные значения зависят от начальной точки и параметров. В таблице приведены характерные порядки.
\end{remark}

\subsubsection{Графики сходимости}

\begin{figure}[H]
    \centering
    \includegraphics[width=\textwidth]{img/convergence.png}
    \caption{Сходимость $|f(\mathbf{x}_k) - f(\mathbf{x}^*)|$ по итерациям}
    \label{fig:convergence}
\end{figure}

На рис.~\ref{fig:convergence} видно:
\begin{itemize}
    \item Метод Ньютона демонстрирует наиболее быструю сходимость на всех функциях.
    \item На хорошо обусловленной квадратичной функции все методы сходятся быстро.
    \item На плохо обусловленной функции градиентный спуск значительно замедляется.
    \item На функции Розенброка градиентный спуск сходится очень медленно из-за узкой долины.
\end{itemize}

\subsubsection{Сходимость по времени}

\begin{figure}[H]
    \centering
    \includegraphics[width=\textwidth]{img/convergence_time.png}
    \caption{Сходимость $|f(\mathbf{x}_k) - f(\mathbf{x}^*)|$ по времени}
    \label{fig:convergence_time}
\end{figure}

\subsubsection{Норма градиента}

\begin{figure}[H]
    \centering
    \includegraphics[width=\textwidth]{img/grad_norm.png}
    \caption{Норма градиента $\|\nabla f(\mathbf{x}_k)\|$ по итерациям (GD и Newton)}
    \label{fig:grad_norm}
\end{figure}

\subsection{Траектории спуска по линиям уровня}

Для визуализации берётся двумерное сечение: фиксируются все координаты в $\mathbf{x}^*$, кроме первых двух ($x_1$, $x_2$).

\begin{figure}[H]
    \centering
    \includegraphics[width=\textwidth]{img/trajectories.png}
    \caption{Траектории методов на линиях уровня (сечение по $x_1$, $x_2$)}
    \label{fig:trajectories}
\end{figure}

На рис.~\ref{fig:trajectories} наглядно видно:
\begin{itemize}
    \item Метод Ньютона идёт напрямую к минимуму.
    \item Градиентный спуск на плохо обусловленной задаче совершает <<зигзаги>> вдоль вытянутых линий уровня.
    \item Нелдер--Мид движется менее предсказуемо, но устойчиво приближается к минимуму.
\end{itemize}

\subsection{Исследование динамики $\varepsilon$ конструктивных чисел}

Ключевое исследование данной работы --- поведение погрешности $\varepsilon$ конструктивных чисел при итеративной оптимизации.

Эксперименты проведены при трёх начальных $\varepsilon_0 \in \{10^{-2}, 10^{-5}, 10^{-8}\}$.

\subsubsection{Общая динамика $\varepsilon$}

\begin{figure}[H]
    \centering
    \includegraphics[width=\textwidth]{img/epsilon_dynamics.png}
    \caption{Динамика среднего $\varepsilon$ по итерациям для разных функций, методов и начальных $\varepsilon_0$}
    \label{fig:eps_dynamics}
\end{figure}

\subsubsection{Покомпонентная динамика $\varepsilon$}

\begin{figure}[H]
    \centering
    \includegraphics[width=\textwidth]{img/epsilon_per_component.png}
    \caption{Покомпонентная динамика $\varepsilon$ для градиентного спуска ($\varepsilon_0 = 10^{-5}$)}
    \label{fig:eps_per_component}
\end{figure}

\subsubsection{Анализ роста $\varepsilon$}

Из результатов экспериментов следуют ключевые наблюдения:

\begin{enumerate}
    \item \textbf{Экспоненциальный рост}: $\varepsilon$ растёт экспоненциально с числом итераций. Каждая арифметическая операция расширяет интервал, и при итеративных вычислениях этот эффект накапливается.

    \item \textbf{Влияние числа обусловленности}: на плохо обусловленной функции ($\kappa = 100$) $\varepsilon$ растёт значительно быстрее, чем на хорошо обусловленной ($\kappa \approx 1$), что связано с большими коэффициентами матрицы $A$.

    \item \textbf{Преимущество Ньютона}: метод Ньютона накапливает наименьшую погрешность, так как требует значительно меньше итераций. Для квадратичных функций --- всего 1 итерация вычислений.

    \item \textbf{Масштабирование}: относительный рост $\varepsilon_{\mathrm{final}} / \varepsilon_0$ практически не зависит от начального $\varepsilon_0$. То есть если начать с $\varepsilon_0 = 10^{-8}$ вместо $10^{-2}$, финальное $\varepsilon$ будет в $10^6$ раз меньше, но коэффициент роста сохранится.

    \item \textbf{Градиентный спуск}: из-за большого числа итераций $\varepsilon$ может расти на десятки порядков, что делает результат бессмысленным с точки зрения интервальных гарантий. Это демонстрирует фундаментальное ограничение наивной интервальной арифметики для итеративных алгоритмов.
\end{enumerate}


% ================================================================
%                         ВЫВОДЫ
% ================================================================
\section{Выводы}

\subsection{По методам оптимизации}

\begin{enumerate}
    \item \textbf{Метод Ньютона} (2-го порядка) --- наиболее эффективен по числу итераций и точности. На квадратичных функциях сходится за 1 итерацию. На функции Розенброка демонстрирует квадратичную скорость сходимости вблизи минимума. Основной недостаток --- необходимость вычисления и обращения гессиана.

    \item \textbf{Градиентный спуск} (1-го порядка) --- скорость сходимости критически зависит от числа обусловленности задачи. При $\kappa \approx 1$ работает хорошо, при $\kappa = 100$ сходится на порядки медленнее. Требует тщательного подбора learning rate.

    \item \textbf{Нелдер--Мид} (0-го порядка) --- не требует производных, что делает его применимым к негладким и шумным функциям. Однако сходится существенно медленнее, особенно при росте размерности.
\end{enumerate}

\subsection{По конструктивным числам}

\begin{enumerate}
    \item Интервальная арифметика обеспечивает гарантированные оценки погрешности, но является пессимистичной --- $\varepsilon$ растёт значительно быстрее, чем реальная ошибка.

    \item Для итеративных алгоритмов рост $\varepsilon$ является экспоненциальным, что ограничивает практическую применимость наивной интервальной арифметики.

    \item Методы, требующие меньше итераций (Ньютон), накапливают меньшую погрешность $\varepsilon$ --- ещё один аргумент в пользу использования методов высших порядков.

    \item Число обусловленности задачи влияет не только на скорость сходимости, но и на скорость роста $\varepsilon$.
\end{enumerate}

\end{document}
